 % ---- ETD Document Class and Useful Packages ---- %
% v1.5.0 released July 7, 2025.
% https://github.com/uchicago-library/uchicago-dissertation
\documentclass[singlespace]{ucetd}

\usepackage{hyperref}
\usepackage[T1]{fontenc}
\usepackage{subcaption,graphicx}
\usepackage{natbib}
\usepackage{mathtools}  % loads amsmath
\usepackage{amssymb}    % loads amsfonts
\usepackage{amsthm}
\usepackage[skip=1em plus 0.5em minus 0.2em]{parskip}
\usepackage{booktabs}


% ----- Chapter micro-abstract (indented both sides, italic) -----
\newcommand{\chapterabstract}[1]{%
  \par\vspace{1em}
  \begingroup
  \leftskip=2em
  \rightskip=2em
  \noindent\itshape
  #1
  \par
  \endgroup
  \vspace{1.5em}
}



%% Use these commands to set biographic information for the title page:
\newcommand{\thesistitle}{FinBERT vs FinGPT Enhancing\\
Malaysian Stock Predictions\\
with News-Driven Insights\\}
\newcommand{\thesisauthor}{Anas Nik Adnan\\}
\department{Computer Science}
\division{Thesis Division}
\degree{MEng Computer Science}
\date{2026}

\title{\thesistitle}
\author{\thesisauthor}

%% Use these commands to set a dedication and epigraph text
\dedication{Dedication Text}
\epigraph{Epigraph Text}

\usepackage{doi}
\usepackage{xurl}
\hypersetup{bookmarksnumbered,
            linktoc=all,
            pdftitle={\thesistitle},
            pdfauthor={\thesisauthor},
            pdfsubject={},                                % Add subject/description
            % pdfkeywords={keyword1, keyword2, keyword3}, % Uncomment and revise keywords
            pdfborder={0 0 0}}
% See https://github.com/k4rtik/uchicago-dissertation/issues/1
\makeatletter
\let\ORG@hyper@linkstart\hyper@linkstart
\protected\def\hyper@linkstart#1#2{%
  \lowercase{\ORG@hyper@linkstart{#1}{#2}}}
\makeatother

\begin{document}
%% Basic setup commands
% If you don't want a title page comment out the next line and uncomment the line after it:
\maketitle
%\omittitle

% These lines can be commented out to disable the copyright/dedication/epigraph pages
\makecopyright

\abstract
This thesis investigates the integration of financial news sentiment into stock market forecasting for the Malaysian equity market. The research compares discriminative and generative language-model approaches for extracting sentiment from Malaysian financial and political news, and studies whether these signals can improve downstream prediction pipelines. The motivation stems from the observation that Malaysian stock prices are often influenced by local and regional news, yet existing predictive systems rarely account for this contextual sentiment. This study therefore addresses two linked challenges: adapting sentiment models trained primarily on Western data to a Malaysian linguistic and market context, and integrating qualitative news information with quantitative time-series market data.

The research consists of three empirical studies.
\begin{enumerate}
    \item \textbf{Sentiment Model Benchmarking}: The first investigation evaluates FinBERT and a family of FinGPT-style generative models on Malaysian financial and political news headlines, assessing zero-shot performance, supervised fine-tuning, and pseudo-label adaptation. The analysis benchmarks sentiment classification accuracy using precision, recall, and macro F1-score, and compares the contextual relevance of sentiment polarity across model families and backbone sizes.
    
    \item \textbf{Sentiment vs. Stock Movement}: The second investigation studies the relationship between daily sentiment scores and Bursa Malaysia returns across multiple time lags ($T$, $T{+}1$, $T{+}2$). Correlation, directional-accuracy, and lag analyses are applied first to a broad market-wide corpus and then to company-filtered corpora in order to distinguish signal dilution from genuine stock-specific association.
    
    \item \textbf{Stock Price Prediction Models}: The final investigation develops predictive models that combine engineered stock features such as moving averages, RSI, MACD, and volatility with benchmark-selected sentiment scores. Fine-tuned FinBERT serves as the primary downstream sentiment signal, while stronger generative baselines are retained as comparative references. LSTM, GRU, and transformer-based models are trained and benchmarked using RMSE, MAE, and directional accuracy metrics.
\end{enumerate}

The completed benchmark and sentiment--return experiments produce three headline findings. First, the fine-tuned FinBERT run2 model remained the strongest sentiment extractor, achieving macro F1-scores of 0.620 on the 2025 Malaysian benchmark and 0.598 on a held-out 2026 test set. Second, larger open instruct backbones improved zero-shot generative performance, with Meta-Llama-3.1-8B-Instruct reaching macro F1 0.458, but supervised LoRA tuning of Qwen-7B and Llama-8B did not surpass either their own zero-shot baselines or FinBERT. Third, broad market-wide sentiment showed negligible association with KLSE index returns, whereas company-filtered sentiment produced materially stronger same-day relationships for Maybank, CIMB, and Tenaga.


The thesis presents following contribution to science:
\begin{enumerate}
    \item \textbf{First Benchmarking of FinBERT vs FinGPT-Style Models on Malaysian News}. This research presents the first comparative analysis between FinBERT and open generative financial LLM workflows for sentiment classification in Malaysian financial and political contexts. It contributes new evidence on model generalization and domain transferability of large language models (LLMs) trained on Western data when applied to emerging markets.
    
    \item \textbf{Multimodal Financial Prediction Framework}. The study introduces an integrated forecasting architecture that fuses sentiment features with quantitative stock indicators, contributing to the advancement of hybrid modelling approaches that bridge natural language understanding and financial time-series prediction.
    
    \item \textbf{Systematic Evaluation of Open Generative Backbones in an Emerging Market Setting}. By comparing TinyLlama-, Qwen-, and Llama-based generative sentiment models against FinBERT on Malaysian news, this research expands evidence on how open Transformer backbones transfer to financial NLP tasks in emerging markets.
    
    \item \textbf{Localized Sentiment Annotation and Analysis}. The creation of a manually annotated Malaysian financial news dataset contributes a valuable resource for future research in regional natural language processing and financial text analytics.
\end{enumerate}

The investigation describes the following impact on industry:
\begin{enumerate}
    \item \textbf{Enhanced Trading and Investment Decisions}.
    Sentiment-augmented prediction models allow local traders and asset managers to react more quickly to market-moving news, improving buy/sell decision accuracy.
    
    \item \textbf{Real-time Market Monitoring Tools}.
    The research enables the development of lightweight dashboards that summarise Malaysian news sentiment, supporting analysts and fintech platforms in daily market assessment.
    
    \item \textbf{Lowering AI Adoption Barriers for Malaysian Fintech}.
    The proposed multimodal framework provides a practical, low-cost template for financial institutions and startups to integrate NLP-based forecasting into their systems.
    
    \item \textbf{Strengthening Data-driven Innovation in Malaysian Capital Markets}.
    By demonstrating the value of LLM-driven sentiment analysis, the work encourages broader adoption of AI analytics, improving market transparency and investor confidence.
\end{enumerate}

\begin{center}
See the project code on
\href{https://github.com/anasamri12/finbert-vs-fingpt.git}{GitHub Repository}.
\end{center}


% Enter Abstract here


\acknowledgments
% Enter Acknowledgements here


\makededication
\makeepigraph


%% Make the various tables of contents
\tableofcontents
\listoffigures
\listoftables



\mainmatter
% Main body of text follows




\chapter{Introduction}


\chapterabstract{%
This chapter introduces the research problem of enhancing financial forecasting in the Malaysian stock market (KLSE) by integrating unstructured news sentiment with deep learning models. It outlines the motivation for moving beyond traditional technical analysis toward Large Language Models (LLMs), comparing the discriminative FinBERT architecture against a family of FinGPT-style generative backbones. The chapter details the research objectives and provides a high-level overview of the three core experiments: sentiment model benchmarking, statistical correlation analysis, and hybrid stock price prediction. Finally, it highlights the study's contributions to science---specifically the evaluation of open Transformer-based sentiment models in an emerging-market setting---and concludes with an outline of the thesis structure.
}

\section{Motivation}

The motivation for this research arises from the inherent complexity of financial markets, in which asset prices are driven not only by fundamental economic indicators but also by the volatile psychology of market participants. Traditional forecasting approaches, including technical analysis and econometric models such as Auto-Regressive Integrated Moving Average (ARIMA), primarily rely on historical price and volume information. However, a growing body of literature demonstrates that such quantitative indicators frequently fail to capture non-linear market shocks triggered by qualitative events, including regulatory changes, political instability, and corporate announcements. Consequently, there is a strong motivation to incorporate unstructured textual data---particularly financial news---into predictive pipelines to improve forecasting performance.

Although sentiment analysis has long been explored as a tool for stock market prediction, the field is currently undergoing a significant technological paradigm shift that remains under-examined in emerging markets. Early sentiment-based approaches relied on lexicon-driven methods, such as the Loughran--McDonald dictionary, which employ bag-of-words representations that ignore semantic context and often misclassify domain-specific financial terminology. The emergence of discriminative deep learning models, most notably FinBERT, marked a substantial advancement by leveraging transfer learning to capture bidirectional contextual information from financial texts. Despite consistently outperforming traditional classifiers, FinBERT remains constrained by its encoder-only architecture: it functions as a static classifier and lacks the ability to perform explicit reasoning or to interpret complex numerical relationships embedded within financial disclosures.

These limitations motivate a shift toward Generative Large Language Models (LLMs). Recent developments, such as FinGPT, employ instruction tuning to reformulate sentiment analysis from a classification task into a generative reasoning task. This enables models to exhibit enhanced numerical sensitivity and contextual inference capabilities that are not readily achievable with encoder-only architectures like FinBERT. However, widely publicised financial LLMs, such as BloombergGPT, remain proprietary and computationally prohibitive for academic research, creating significant barriers to reproducibility and accessibility. FinGPT addresses this challenge by providing an open-source framework that utilises Low-Rank Adaptation (LoRA) to enable efficient fine-tuning of large models at substantially reduced computational cost.

Critically, a major research gap exists in the application of these advanced architectures to the Malaysian stock market (KLSE). The majority of financial Natural Language Processing research is concentrated on highly liquid and information-rich markets in the United States and China. In contrast, the KLSE exhibits distinct microstructural characteristics, including heightened sensitivity to local political developments and unique investor behaviour patterns that may not be adequately captured by globally trained models without domain-specific adaptation. While recent Malaysian studies have begun to apply machine learning techniques to stock prediction, they typically rely on lexicon-based sentiment analysis or conventional Long Short-Term Memory (LSTM) models, without exploiting the reasoning capabilities of modern generative LLMs.

Accordingly, this research is motivated by the need to empirically evaluate whether open generative financial LLMs can match or exceed the static sentiment classification provided by FinBERT within the Malaysian market context. By integrating the strongest benchmarked sentiment signals into LSTM-based time-series forecasting frameworks, this study seeks to determine whether the additional computational complexity associated with Generative AI translates into meaningful predictive gains for investors operating in emerging markets.

\section{Research Objectives}

Based on the problem statement and the research gaps identified within the Malaysian market context, this thesis proposes a set of research objectives designed to empirically evaluate the trade-off between the computational efficiency of discriminative models and the generative reasoning capabilities of Large Language Models (LLMs) in financial forecasting.

The primary aim of this research is to develop a robust framework that integrates advanced Natural Language Processing (NLP) techniques with deep learning models to predict stock price movements in Bursa Malaysia (KLSE). To achieve this aim, the study addresses the following specific objectives:

\begin{enumerate}
    \item \textbf{Benchmark sentiment classification performance of discriminative versus generative models on Malaysian financial news.}  
    This objective involves fine-tuning FinBERT, a discriminative encoder-based model, and benchmarking a family of generative instruction-tuned models using a curated dataset of Malaysian financial news. The generative side begins with a TinyLlama-based FinGPT baseline and is extended to larger Qwen and Llama instruct backbones in order to determine which architectural paradigm and backbone scale more effectively capture local linguistic nuances, financial jargon, and the pronounced neutral-class imbalance characteristic of Malaysian corporate reporting.

    \item \textbf{Quantify the relationship between extracted sentiment signals and KLSE stock returns.}  
    This objective seeks to statistically analyse whether sentiment scores produced by the strongest benchmarked models are associated with realised price movement in Bursa Malaysia. The analysis first measures broad corpus correlations against the KLSE index and major stocks, and then tests whether company-specific news filtering yields stronger stock-level relationships.

    \item \textbf{Develop and evaluate hybrid deep learning models for time-series forecasting.}  
    This objective focuses on constructing hybrid forecasting pipelines that integrate unstructured sentiment information with structured historical price data (Open, High, Low, Close, and Volume). By incorporating sentiment scores as auxiliary features within Long Short-Term Memory (LSTM) networks, the study assesses whether narrative-driven signals significantly reduce prediction error, as measured by Root Mean Square Error (RMSE) and Mean Absolute Percentage Error (MAPE), relative to price-only baselines.

    \item \textbf{Evaluate the computational trade-off between model complexity and predictive performance.}  
    The final objective examines whether the increased computational cost and inference latency associated with generative models such as FinGPT yield a statistically significant improvement in forecasting accuracy compared to lightweight discriminative models like FinBERT. This analysis is critical for assessing the practical feasibility of deploying Generative AI in real-time or near--real-time trading and decision-support systems within emerging market environments.
\end{enumerate}

\section{Research Experiments}

This thesis adopts a progressive empirical design that begins with an evaluation of the linguistic capabilities of domain-specific sentiment models, advances to a statistical analysis of the relationship between sentiment signals and market returns, and culminates in an out-of-sample forecasting study. The research is structured into three interconnected empirical experiments, each addressing a distinct aspect of the proposed research objectives.

\begin{enumerate}
    \item \textbf{Benchmarking Discriminative versus Generative Sentiment Models on Malaysian Financial Corpora}  
    The first experiment seeks to empirically characterise the trade-off between \emph{classification accuracy} and \emph{contextual reasoning} within the Malaysian financial domain. FinBERT, a discriminative encoder-based model, has established strong performance in financial sentiment classification through supervised fine-tuning; however, it may struggle with implied or context-dependent sentiment commonly present in complex financial narratives. In contrast, generative FinGPT-style models leverage instruction tuning to reformulate sentiment analysis as a text generation task.  
    This experiment benchmarks both paradigms using a curated dataset of Malaysian financial and political news sourced from outlets such as \emph{The Star} and \emph{New Straits Times}, and extends the generative comparison from an initial TinyLlama baseline to larger Qwen and Llama instruct backbones. Zero-shot prompting, supervised fine-tuning, and pseudo-label adaptation are all evaluated, with particular emphasis on neutral-class imbalance and local economic terminology that generic models frequently misclassify.

    \item \textbf{Statistical Analysis of Market Sensitivity to Sentiment Shocks}  
    The second experiment examines whether sentiment signals extracted from Malaysian financial news are associated with daily returns in Bursa Malaysia. Moving beyond classification metrics, the study first evaluates broad market-wide sentiment against the KLSE index and selected large-cap stocks, and then repeats the analysis on company-filtered corpora to test whether article relevance materially changes the signal. Pearson and Spearman correlation, directional accuracy, and lag analysis are used to distinguish same-day association from genuine next-day predictive value.

    \item \textbf{Hybrid Forecasting with Benchmark-Selected Sentiment Signals}  
    The final experiment integrates insights from the preceding studies to develop and evaluate hybrid forecasting architectures. Traditional econometric models and standalone Long Short-Term Memory (LSTM) networks often struggle to capture non-linear dependencies induced by qualitative market events. In this experiment, sentiment representations are fused with historical OHLCV (Open, High, Low, Close, Volume) data to form enriched input sequences for time-series prediction. A sliding window validation strategy is employed to prevent look-ahead bias. The predictive performance of a baseline LSTM model is compared against hybrid models incorporating sentiment features from the strongest Chapter~3 sentiment extractors, with fine-tuned FinBERT serving as the primary downstream signal and stronger generative baselines retained as comparative references where useful. Model performance is assessed using Root Mean Square Error (RMSE) and Mean Absolute Percentage Error (MAPE), with the primary objective of determining whether localized sentiment improves forecasting accuracy over lightweight price-only baselines in the Malaysian market context.
\end{enumerate}


\section{Scientific Contributions}


This research contributes to the existing literature in several significant ways:

\begin{enumerate}
    \item \textbf{First Comparative Benchmarking of FinBERT and FinGPT-Style Models in the Malaysian Market}  
    To the best of the author's knowledge, this study presents the first empirical comparison between discriminative (encoder-only) and generative (decoder-only) financial language models within the context of Bursa Malaysia (KLSE). While FinBERT has been extensively evaluated on developed markets such as the United States and China, open generative financial LLM workflows have rarely been benchmarked in Malaysian news settings. This thesis addresses that gap by comparing FinBERT against TinyLlama-, Qwen-, and Llama-based generative backbones under zero-shot, supervised, and pseudo-label adaptation settings.

    \item \textbf{Systematic Evaluation of Open Generative Backbones and Adaptation Strategies}  
    This research advances the methodological understanding of how open generative models transfer to financial sentiment analysis in an emerging market. Rather than treating FinGPT as a single fixed model, the thesis compares multiple backbone scales and families, and evaluates zero-shot prompting, supervised LoRA fine-tuning, and pseudo-label adaptation. The resulting benchmark provides new evidence on how backbone capacity and adaptation strategy affect generative financial NLP performance on Malaysian news.

    \item \textbf{Development of a Hybrid LLM--LSTM Forecasting Framework}  
    This thesis contributes to financial time-series forecasting by developing a hybrid predictive architecture that integrates unstructured sentiment signals with structured market data (Open, High, Low, Close, and Volume). In contrast to prior studies that rely on lexicon-based sentiment measures or standalone LSTM models, this research directly incorporates Transformer-derived sentiment representations into the LSTM forecasting pipeline. The proposed framework empirically evaluates whether narrative-driven sentiment information leads to meaningful reductions in forecasting error, as measured by RMSE and MAPE, relative to unimodal price-based baselines.

    \item \textbf{Creation of a Labelled Malaysian Financial News Dataset}  
    A persistent challenge in applying Natural Language Processing to emerging markets is the scarcity of high-quality annotated data. This research contributes a curated and sentiment-labelled dataset of Malaysian financial news, addressing the low-resource limitation frequently cited in the literature. This dataset provides a valuable benchmark for future research on financial NLP in Southeast Asia and supports reproducibility in emerging-market studies.

    \item \textbf{Assessment of Computational Trade-offs in Financial Artificial Intelligence}  
    Finally, this study offers a practical evaluation of the trade-off between model complexity and predictive performance in financial applications. By contrasting the lightweight and efficient inference of FinBERT with the higher computational cost and latency associated with larger generative backbones, the research provides actionable insights into whether the deployment of such models is justified for real-time or near--real-time trading environments in markets with lower liquidity than major global exchanges.
\end{enumerate}


\section{Thesis Structure}
The structure of this thesis is organised as follows:
\begin{itemize}
    \item Chapter 2 - \textbf{Background and Literature Review}. The relevant literature and key concepts in the areas of this research are reviewed to introduce the problem and environment of this thesis. The chapter starts by presenting the application domain, specifically the Malaysian stock market (KLSE) and the role of news sentiment in financial volatility. It reviews the evolution of sentiment analysis from lexicon-based approaches to Large Language Models (LLMs) like FinBERT and FinGPT. The chapter concludes with a methodological review of time-series forecasting (LSTM, GRU) and the integration of unstructured text data into financial prediction models.

    \item Chapter 3 - \textbf{Sentiment Model Benchmarking}. This chapter provides an empirical benchmarking analysis to compare the capabilities of FinBERT and FinGPT-style generative models in extracting sentiment from Malaysian financial and political news. The study covers zero-shot prompting, supervised fine-tuning, and pseudo-label adaptation on a curated localized benchmark, and identifies the strongest sentiment model for later integration into predictive pipelines.

    \item Chapter 4 - \textbf{Sentiment vs. Stock Movement}. This second study explores the relationship between the daily sentiment scores derived in Chapter 3 and the actual market movements of top Malaysian companies. The analysis contrasts broad market-wide sentiment against company-filtered sentiment, and uses lag analysis ($T$, $T{+}1$, $T{+}2$), Pearson/Spearman correlation, and directional accuracy to test whether sentiment behaves as a same-day association signal or a genuine leading indicator.

    \item Chapter 5 - \textbf{Stock Price Prediction Models}. This study integrates the findings of the previous experiments to build a hybrid predictive framework. It proposes time-series prediction models that fuse historical stock data with benchmark-selected sentiment features. The chapter compares baseline models (LSTM/GRU with price features only) against sentiment-enhanced models, using the strongest Chapter~3 sentiment signal as the primary downstream feature and retaining generative baselines as comparative references where appropriate.

    \item Chapter 6 - \textbf{Conclusion and Future Work}. The final chapter provides an overall conclusion of this research with a summary of the key findings regarding the applicability of LLMs in emerging markets. The thesis ends with recommendations for future work, such as the development of real-time market monitoring tools and the expansion of the framework to other ASEAN economies.
\end{itemize}

% Introductory stuff




\chapter{Background and Literature Review}

\chapterabstract{This chapter critically examines the theoretical foundations of financial sentiment analysis within the specific context of the Malaysian stock market (KLSE). It traces the methodological evolution of Natural Language Processing (NLP) in finance, highlighting the transition from static lexicon-based approaches to dynamic Large Language Models (LLMs), specifically contrasting the architectural capabilities of encoder-only models (FinBERT) against generative decoder models (FinGPT). Furthermore, the chapter evaluates the existing literature on time-series forecasting, identifying a significant gap in hybrid modelling where unstructured, localized news data is integrated with Deep Learning algorithms (LSTM, GRU) to predict market volatility in emerging economies.}

\section{Application Domain}
\subsection{Application Domain: The Malaysian Stock Market (KLSE)}

The application domain of this thesis is Bursa Malaysia, specifically the Kuala Lumpur Composite Index (KLCI), which serves as the primary barometer for the Malaysian equity market. As one of the largest bourses in the Association of Southeast Asian Nations (ASEAN), the KLSE exhibits unique microstructural characteristics that distinguish it from the developed markets (e.g.\ NYSE, NASDAQ) typically used to train financial Large Language Models (LLMs). Unlike mature markets that are driven largely by institutional and algorithmic trading, the Malaysian market is characterised by significant volatility influenced by local socio-political events and retail investor sentiment.

\subsubsection{Sensitivity to News and Events}

The literature establishes that the KLSE is highly sensitive to unstructured information flow, where market volatility often correlates more strongly with news sentiment than with fundamental macroeconomic indicators alone. Historical analyses reveal that major price fluctuations in the KLSE are frequently event-driven. For instance, the market experienced severe volatility during the 1997 Asian Financial Crisis, the 2008 Global Financial Crisis, and major domestic political events such as the 14th Malaysian General Election (GE14). These events demonstrate that stock prices in Malaysia are not merely reflections of quantitative financial ratios but are heavily influenced by public mood and news coverage surrounding such crises. Consequently, traditional forecasting models that rely solely on historical price data (e.g.\ OHLCV) or technical indicators often fail to capture the directional shifts caused by these qualitative factors.

\subsubsection{The Gap in Localised Sentiment Analysis}

While sentiment analysis has been widely applied to Western financial markets, the literature reveals a significant gap in the Malaysian context. Most existing studies on the KLSE rely on global sentiment datasets or generic social media sources (e.g.\ Twitter/X), which do not adequately reflect regional investor behaviour, local language nuances, or the specific economic context of Malaysia. Research by Tan \textit{et al.}\ (2025) argues that sentiment data must be \emph{localised}---extracted from domestic platforms such as KLSE Screener or local news portals---to accurately model the Malaysian financial ecosystem.

Furthermore, standard sentiment lexicons (e.g.\ VADER, Loughran--McDonald) often misclassify financial news in the Malaysian context due to linguistic diversity and domain-specific jargon commonly used in local reporting. This limitation necessitates the use of advanced deep learning and Transformer-based models (such as FinBERT and FinGPT), which can be fine-tuned to capture the semantic context of Malaysian financial news. In contrast, static dictionary-based approaches have demonstrated limited accuracy in this domain.

\subsubsection{Justification for LLM Integration}

The integration of Artificial Intelligence (AI) into Malaysian stock prediction remains an under-explored research area. Although recent studies have employed Long Short-Term Memory (LSTM) and Convolutional Neural Network (CNN) architectures for KLSE forecasting, there is a scarcity of work utilising Generative AI and Large Language Models (LLMs) to process unstructured financial text in this market. By focusing on the KLSE, this thesis addresses the critical need to evaluate whether LLMs pre-trained on Western financial data (such as FinBERT) can effectively transfer to emerging markets, or whether generative approaches (such as FinGPT) provide superior adaptability for capturing local market dynamics.



\subsection{The Evolution of Financial Sentiment Analysis}

The methodology for extracting sentiment from financial text has undergone a significant paradigm shift over the past decade, transitioning from static, dictionary-based approaches to dynamic, context-aware Large Language Models (LLMs). This evolution reflects the growing need to move beyond simple word-frequency analysis toward modelling the complex semantic structures, contextual dependencies, and numerical reasoning inherent in financial reporting.

\subsubsection{Lexicon-Based Approaches}

Early research in financial sentiment analysis relied predominantly on lexicon (dictionary-based) approaches. These methods compute sentiment by matching words in a document against predefined lists of positive and negative terms. A foundational contribution in this area was the development of the Loughran--McDonald (LM) financial dictionary. Loughran and McDonald (2011) demonstrated that general-purpose sentiment lexicons (such as Harvard IV-4) are ill-suited for financial applications; for example, words such as \emph{``liability''} or \emph{``tax''} are often classified as negative in general language but are neutral descriptors in financial statements.

While lexicon-based methods such as VADER and the LM dictionary remain in use due to their computational efficiency and interpretability, they suffer from a fundamental \emph{bag-of-words} limitation. These approaches treat words independently and fail to capture semantic dependencies or contextual meaning. As a result, they frequently misclassify financial news where sentiment is implied rather than explicitly expressed through emotional keywords. In the Malaysian context, researchers have applied dictionary-based sentiment measures alongside algorithms such as Hybrid Na\"ive Bayes to correlate news sentiment with local stock price movements; however, their performance remains limited by these structural constraints.

\subsubsection{Machine Learning and Deep Learning Approaches}

To overcome the shortcomings of lexicon-based methods, researchers introduced supervised Machine Learning (ML) algorithms such as Support Vector Machines (SVM) and Random Forest classifiers. In these frameworks, sentiment analysis is formulated as a classification task using features such as N-grams to capture limited contextual information. Although SVM-based models improved classification accuracy relative to static dictionaries, they continued to struggle with the non-linear and long-range dependencies present in natural language.

The emergence of Deep Learning marked a further methodological advancement. Word embedding techniques (e.g.\ Word2Vec and GloVe) enabled words to be represented as dense vectors, allowing models to capture semantic similarity between terms. Recurrent Neural Networks (RNNs), particularly Long Short-Term Memory (LSTM) networks, became widely adopted for sentiment analysis due to their ability to process sequential data and mitigate the vanishing gradient problem. Empirical studies consistently demonstrate that LSTM-based models outperform traditional ML classifiers in predicting stock market movements from news text, owing to their improved handling of temporal and contextual dependencies.

\subsubsection{The Transformer Revolution: FinBERT}

The introduction of the Transformer architecture and Bidirectional Encoder Representations from Transformers (BERT) represented a major breakthrough in Natural Language Processing (NLP). By leveraging self-attention mechanisms, BERT models learn contextual representations by considering both preceding and following tokens simultaneously. However, generic BERT models pre-trained on corpora such as Wikipedia lack sufficient understanding of financial terminology and discourse.

FinBERT addresses this limitation through domain-specific pre-training on large-scale financial corpora, including corporate filings (e.g.\ 10-K reports), earnings call transcripts, and financial news articles. Architecturally, FinBERT is a discriminative, encoder-only model that adds a dense classification layer on top of the transformer output to assign sentiment labels (Positive, Negative, or Neutral). Empirical benchmarks indicate that FinBERT significantly outperforms generic BERT and LSTM-based models in financial sentiment classification due to its ability to capture domain-specific language and financial semantics.

\subsubsection{Generative AI and Large Language Models}

The most recent stage in the evolution of financial sentiment analysis is the emergence of Generative Artificial Intelligence and decoder-only Large Language Models (LLMs), such as GPT-4, BloombergGPT, and FinGPT. Unlike discriminative models that map text to predefined labels, generative LLMs can produce free-form outputs, perform reasoning, and follow complex instructions.

While proprietary models such as BloombergGPT are trained on extensive internal financial datasets, open-source frameworks such as FinGPT have expanded accessibility through parameter-efficient fine-tuning techniques, including Low-Rank Adaptation (LoRA). These methods enable large pre-trained LLMs (e.g.\ LLaMA-based architectures) to be adapted to financial tasks using comparatively modest computational resources.

A key advantage of generative LLMs over models such as FinBERT lies in their numerical sensitivity. Financial news frequently contains quantitative statements (e.g.\ \emph{``revenue exceeded estimates by 5\%''}), where sentiment depends on numerical interpretation rather than explicit emotional language. FinBERT often struggles in such cases, as numerical implications are weakly represented in its classification objective. In contrast, instruction-tuned generative models such as FinGPT demonstrate superior numerical reasoning capabilities, allowing them to infer sentiment from implicit quantitative and contextual cues. This capability is particularly important in the Malaysian market, where financial reporting commonly intertwines numerical disclosures with political and macroeconomic narratives.

\subsubsection{Research Questions Motivating This Thesis}

Based on the reviewed literature, this thesis is motivated by the following research questions:

\begin{itemize}
    \item How do FinBERT and FinGPT differ in their sensitivity to numerical information in financial news?
    \item What unique challenges does the KLSE market present for Large Language Models?
    \item Why might instruction-tuned generative models outperform classification-based approaches for financial sentiment analysis?
\end{itemize}



\subsection{Large Language Models: Discriminative vs.\ Generative}

The advent of the Transformer architecture has bifurcated the landscape of Natural Language Processing (NLP) into two dominant paradigms: discriminative models based on encoder-only architectures, and generative models based on decoder-only architectures. This section critically evaluates these approaches within the financial domain, highlighting the theoretical trade-offs between computational efficiency and contextual reasoning that motivate the core research questions of this thesis.

\subsubsection{Discriminative Models: The Encoder-Only Paradigm (FinBERT)}

Discriminative models employ the Transformer encoder to produce contextualised representations of input text, primarily for downstream classification tasks. The foundational model in this category is Bidirectional Encoder Representations from Transformers (BERT), which is trained using a masked language modelling (MLM) objective to learn deep bidirectional semantic representations.

In financial applications, generic BERT models pre-trained on general-purpose corpora (e.g.\ Wikipedia) often fail to capture the semantic nuances and domain-specific terminology prevalent in financial discourse. To address this limitation, FinBERT was introduced as a domain-adapted variant, pre-trained on large-scale financial corpora including corporate filings (e.g.\ 10-K and 10-Q reports), earnings call transcripts, and financial news articles. Architecturally, FinBERT formulates sentiment analysis as a discriminative classification task: a sequence of tokens is processed by the encoder, and the final hidden representation of the special \texttt{[CLS]} token is passed through a dense classification head to output probabilities over sentiment classes (Positive, Negative, Neutral).

Empirical studies demonstrate that FinBERT consistently achieves state-of-the-art performance on benchmark datasets such as the Financial PhraseBank, outperforming generic BERT and LSTM-based models by substantial margins. However, as an encoder-only architecture, FinBERT is inherently limited by its static design. It requires task-specific classification heads and lacks the ability to perform generative reasoning or explain the underlying rationale behind a sentiment prediction. Moreover, discriminative models often exhibit weak numerical sensitivity, struggling to interpret the significance and directionality of numerical information embedded within financial text. This limitation is particularly problematic in market analysis, where sentiment frequently depends on quantitative changes rather than explicit affective language.

\subsubsection{Generative Models: The Decoder-Only Paradigm (FinGPT)}

The recent paradigm shift toward Generative Artificial Intelligence is driven by decoder-only Transformer architectures (e.g.\ GPT-3, LLaMA), trained using an autoregressive next-token prediction objective. Unlike discriminative models, generative Large Language Models (LLMs) exhibit emergent capabilities such as few-shot learning, instruction following, and chain-of-thought reasoning, enabling them to perform a wide range of tasks without extensive task-specific fine-tuning.

BloombergGPT represents the first large-scale application of this paradigm to the financial domain, comprising a 50-billion-parameter model trained on a mixed corpus of approximately 363 billion financial tokens and 345 billion general-domain tokens. While BloombergGPT demonstrated that combining domain-specific and general knowledge yields strong performance across financial NLP tasks, its proprietary nature and extreme computational requirements---estimated at over one million GPU hours---render it inaccessible for most academic research.

These constraints motivated the development of FinGPT, an open-source framework designed to democratise access to financial LLMs. FinGPT adopts a data-centric methodology and leverages parameter-efficient fine-tuning techniques such as Low-Rank Adaptation (LoRA). By freezing the majority of pre-trained model weights and introducing trainable low-rank matrices, FinGPT reduces the number of trainable parameters from billions to millions, significantly lowering computational cost while preserving model expressiveness.

Crucially, generative models such as FinGPT differ from FinBERT in their contextual understanding and numerical sensitivity. Through instruction tuning, sentiment analysis is reframed from a fixed classification task into a generative reasoning task (e.g.\ \emph{``Determine the sentiment of the following news article''}). This formulation allows the model to integrate broader contextual knowledge and infer sentiment from implicit cues, particularly in complex narratives involving numerical changes, expectations, or macroeconomic implications.

\subsubsection{Synthesis and Model Selection}

The literature highlights a fundamental trade-off between discriminative and generative approaches. FinBERT offers high computational efficiency and strong performance on well-defined sentiment classification benchmarks, making it suitable for large-scale, low-latency processing. However, its decision-making process remains opaque, and its limited reasoning and numerical interpretation capabilities constrain its applicability in complex market environments.

In contrast, FinGPT provides superior adaptability, reasoning ability, and interpretability through natural language generation and prompt-based interaction. These advantages come at the cost of increased inference latency and computational overhead relative to lightweight discriminative models. 

This thesis empirically investigates this trade-off within the Malaysian stock market context. While FinBERT has demonstrated strong performance on US and European datasets, its effectiveness in capturing the linguistic diversity, numerical subtleties, and event-driven dynamics of the KLSE remains under-explored. Conversely, generative models such as FinGPT may offer enhanced capability to reason over local market narratives through few-shot learning and instruction-based inference.



\subsection{Time-Series Forecasting with Hybrid Models}

The transition from traditional econometric models to deep learning architectures represents a major advancement in financial time-series forecasting. Linear models such as Auto-Regressive Integrated Moving Average (ARIMA) have historically served as benchmarks for stationarity-based forecasting; however, they often fail to capture the non-linear volatility patterns and long-term dependencies inherent in stock market data. As a result, the literature has increasingly shifted toward Recurrent Neural Networks (RNNs), and in particular Long Short-Term Memory (LSTM) networks, as the standard architecture for modelling sequential financial data.

\subsubsection{Long Short-Term Memory (LSTM) Networks}

Standard RNNs suffer from the \emph{vanishing gradient} problem, which limits their ability to learn dependencies over long time horizons---a critical limitation when analysing financial trends that may persist over weeks or months. Hochreiter and Schmidhuber (1997) addressed this issue through the introduction of the LSTM architecture, which incorporates a gating mechanism comprising input, forget, and output gates. These gates regulate the flow of information through the network, enabling the model to selectively retain relevant historical information while discarding noise. As a result, LSTM networks maintain an internal memory cell that preserves long-range contextual information across time steps.

Empirical studies consistently demonstrate the superiority of LSTM networks over traditional machine learning models such as Support Vector Machines (SVM) and Multilayer Perceptrons (MLP) for financial time-series forecasting. For example, Halder (2024) reported that LSTM models achieved a 20.23\% improvement in Mean Absolute Percentage Error (MAPE) compared to MLP models when forecasting NASDAQ-100 indices. Similarly, in the Malaysian context, Sidek \textit{et al.}\ (2023) found that LSTM-based architectures outperformed Decision Tree models in minimizing Root Mean Square Error (RMSE) for Bursa Malaysia constituent stocks.

\subsubsection{The Hybrid Approach: FinBERT--LSTM}

A key limitation of applying LSTM models exclusively to price-based data (e.g.\ OHLCV) is the exclusion of qualitative market drivers. To overcome this limitation, recent studies have proposed hybrid architectures that integrate structured numerical features with unstructured sentiment information extracted using Natural Language Processing (NLP) models. This approach is theoretically motivated by the observation that while historical prices capture underlying trends, news sentiment reflects market shocks, investor psychology, and exogenous events that often trigger trend reversals.

The FinBERT--LSTM pipeline represents a state-of-the-art hybrid architecture in this domain. Within this framework, FinBERT acts as a feature extractor, transforming financial text into sentiment representations (e.g.\ classification probabilities or embedding vectors). These sentiment features are then incorporated as auxiliary inputs to the LSTM network alongside traditional technical indicators. Gu \textit{et al.}\ (2024) demonstrated that a FinBERT--LSTM model achieved an accuracy of 0.955 in stock trend prediction, outperforming standalone LSTM models (0.928) and conventional Deep Neural Networks (0.78). 

Further extending this approach, Soliman \textit{et al.}\ (2025) compared LSTM-based architectures with probabilistic forecasting models such as DeepAR. While probabilistic models provided advantages in uncertainty quantification, the integration of FinBERT-derived sentiment features consistently reduced forecasting errors, including RMSE and MAPE, across both deterministic and probabilistic frameworks.



\subsection{Synthesis and Research Gap}

The reviewed literature reveals an increasingly sophisticated financial forecasting landscape in which Deep Learning models and Large Language Models (LLMs) are converging. Despite these methodological advances, several critical research gaps remain, particularly in the context of emerging markets such as Malaysia.

\subsubsection{The Geographic Gap: Under-representation of the KLSE}

The majority of advanced sentiment analysis and financial forecasting research focuses on developed markets, most notably the United States (NASDAQ, NYSE) and China. In contrast, rigorous applications of Transformer-based LLMs to the Bursa Malaysia (KLSE) remain scarce. While recent local studies have explored machine learning approaches for KLSE prediction, they often rely on legacy lexicon-based sentiment techniques (e.g.\ VADER, Loughran--McDonald) or conventional classifiers such as Na\"ive Bayes and SVM. 

Although Tan \textit{et al.}\ (2025) introduced localised sentiment data sourced from platforms such as \emph{KLSE Screener}, their analysis focused primarily on ensemble tree-based models (e.g.\ XGBoost) rather than evaluating the semantic reasoning capabilities of modern Generative AI models. This thesis directly addresses this gap by applying advanced FinBERT and FinGPT architectures to the Malaysian market, with particular emphasis on their ability to capture local political and economic narratives that are often absent from global datasets.

\subsubsection{The Methodological Gap: Discriminative vs.\ Generative Models for Forecasting}

A further under-explored area in the literature is the direct performance comparison between discriminative models (e.g.\ FinBERT) and generative models (e.g.\ FinGPT) when embedded within a downstream time-series forecasting pipeline. Existing studies typically evaluate these models in isolation: FinBERT is recognised for its efficiency and classification accuracy, while FinGPT is examined primarily for its generative and instruction-following capabilities. There is limited empirical evidence quantifying the trade-off between the computational efficiency of encoder-only models and the enhanced contextual reasoning of decoder-only models when their outputs are used as predictive features in emerging market settings.

\subsubsection{Conclusion}

Existing research consistently demonstrates that hybrid models combining sentiment analysis with time-series forecasting (e.g.\ LLM + LSTM) outperform unimodal approaches based solely on price data. However, the application of Generative Financial LLMs such as FinGPT to the unique microstructure of the Malaysian stock market remains largely unexplored. This thesis aims to bridge this gap by empirically evaluating whether the generative reasoning capabilities of FinGPT yield a statistically significant improvement over the static sentiment classification provided by FinBERT in modelling KLSE stock movement. In doing so, this work contributes a novel empirical perspective to the broader discourse on the role of Generative AI in financial market prediction.


\section{Method Development}
% Intro to chapter one

\subsection{Static vs.\ Dynamic Data Representation}

A fundamental methodological distinction in this research is the treatment of input data as either \emph{static} (cross-sectional) or \emph{dynamic} (longitudinal). This distinction determines both the mathematical transformations applied to the data and the choice of model architecture for each branch of the hybrid prediction pipeline.

\paragraph{Static Representation: Unstructured Text and Sentiment}

In the context of Experiment~1 (Sentiment Benchmarking), financial news headlines are treated as static, cross-sectional data. Although natural language is inherently sequential at the syntactic level, the sentiment extraction process treats each headline as an independent observation occurring at a discrete point in time~$t$, without requiring temporal dependence on observations at $t-1$ for classification.

\begin{itemize}
    \item \textbf{Vectorisation and Embedding:}  
    Raw textual data is unstructured and must be transformed into a numerical representation to be processed by discriminative models such as FinBERT. This is achieved through Transformer-based embedding, where the input token sequence is mapped to contextualised representations. The final hidden state of the special \texttt{[CLS]} token serves as a static, high-dimensional vector summarising the sentiment of the entire headline.

    \item \textbf{Dimensionality Reduction:}  
    Unlike traditional bag-of-words approaches that produce sparse and high-dimensional frequency vectors, Transformer models compress semantic information into dense embeddings. For integration into the downstream forecasting model, this representation is further reduced to a scalar sentiment score or a probability distribution over sentiment classes (Positive, Negative, Neutral). This output constitutes a static sentiment feature associated with a specific trading day, denoted as $S_t$.
\end{itemize}

\paragraph{Dynamic Representation: Financial Time-Series}

In contrast, the stock market data used in Experiments~2 and~3---comprising Open, High, Low, Close, and Volume (OHLCV)---is dynamic and longitudinal in nature. The temporal ordering of observations is essential; financial time-series data is not independent and identically distributed (i.i.d.), but instead exhibits serial dependence, where the state at time~$t$ is a function of prior states $t-1, t-2, \dots, t-n$.

\begin{itemize}
    \item \textbf{Sequential Dependencies:}  
    Financial econometrics assumes that historical price movements and trading volumes encode latent structures such as momentum, volatility clustering, and mean reversion. Capturing these dynamics requires representations that preserve chronological ordering, allowing the model to learn state evolution over time.

    \item \textbf{Sliding Window Approach:}  
    To convert dynamic financial data into a supervised learning format suitable for Long Short-Term Memory (LSTM) networks, this study employs a sliding (rolling) window methodology. The resulting input data is structured as a three-dimensional tensor of shape $(\text{Samples}, \text{Time Steps}, \text{Features})$.
    
    For a window size $w$ (e.g.\ $w=30$ trading days), the input vector at time~$t$ is defined as a sequence:
    \[
        \mathbf{X}_t = \{ \mathbf{x}_{t-w}, \mathbf{x}_{t-w+1}, \dots, \mathbf{x}_{t} \}.
    \]
    This formulation ensures that predictions for $t+1$ are generated using only information available up to time~$t$, thereby preventing look-ahead bias.
\end{itemize}

\paragraph{Fusion of Modalities}

The key methodological contribution of this thesis lies in the integration of static textual sentiment with dynamic financial time-series data. The static sentiment scores $S_t$ generated by FinBERT or FinGPT are fused with the dynamic market variables $P_t$ to form an enriched, multimodal dataset.

\begin{itemize}
    \item \textbf{Temporal Alignment:}  
    Financial news arrives irregularly, whereas stock prices are observed at regular daily intervals. To reconcile this mismatch, sentiment signals are aggregated (e.g.\ via averaging or weighted aggregation) to align with the daily frequency of market data.

    \item \textbf{Augmented Input Vector:}  
    The final input to the predictive model is a dynamic sequence in which each time step contains both numerical market indicators and the corresponding static sentiment feature derived from that day's news:
    \[
        \mathbf{X}_t =
        \big[
        \text{Open}_t,\,
        \text{Close}_t,\,
        \text{Volume}_t,\,
        \text{Sentiment}_t
        \big].
    \]
\end{itemize}

\subsection{Supervised Learning Frameworks}

This research employs two distinct supervised learning paradigms to benchmark the effectiveness of financial sentiment analysis: \emph{Discriminative Fine-Tuning}, applied to FinBERT, and \emph{Generative Instruction Tuning}, applied to FinGPT. Although both approaches aim to learn a mapping from input text $\mathbf{X}$ to sentiment labels $\mathbf{Y}$, they differ fundamentally in their architectural objectives, learning formulations, and parameter update mechanisms.

\paragraph{Discriminative Fine-Tuning (The FinBERT Approach)}

FinBERT is based on a Transformer encoder architecture. Its adaptation to the Malaysian financial domain is achieved through \emph{transfer learning}, whereby knowledge learned from a source domain (large-scale financial corpora such as TRC2-financial) is transferred to a target domain consisting of Malaysian financial news.

\begin{itemize}
    \item \textbf{Mechanism:}  
    The input text is tokenised and passed through the BERT encoder. The final hidden representation of the special classification token \texttt{[CLS]} is used as an aggregate semantic embedding of the entire input sequence. A task-specific dense output layer, implemented as a linear projection followed by a Softmax activation, is added on top of the \texttt{[CLS]} representation to produce a probability distribution over sentiment classes (Positive, Negative, Neutral).

    \item \textbf{Optimisation:}  
    During fine-tuning, model parameters are updated by minimising the categorical Cross-Entropy Loss between the predicted class probabilities and the ground truth sentiment labels from the annotated dataset. To mitigate \emph{catastrophic forgetting}---where pre-trained linguistic knowledge is overwritten during task adaptation---this study adopts discriminative fine-tuning. In this scheme, lower encoder layers (which capture general syntactic and semantic features) are assigned smaller learning rates, while higher layers (which encode task-specific information) are updated more aggressively.
\end{itemize}

\paragraph{Generative Instruction Tuning (The FinGPT Approach)}

In contrast to FinBERT, FinGPT-style models are built upon decoder-only Transformer architectures and are designed for autoregressive text generation. In this thesis, the generative family is instantiated with multiple open backbones, including TinyLlama, Qwen, and Llama instruct models, and adapted to sentiment classification through \emph{instruction tuning}.

\begin{itemize}
    \item \textbf{Task Transformation:}  
    Rather than projecting a fixed-length embedding to a sentiment label, FinGPT reformulates sentiment analysis as a text generation task. The dataset is converted into instruction--response pairs of the form:
    \[
        \text{Human: [Instruction] + [Input]} \;\rightarrow\;
        \text{Assistant: [Output]}.
    \]
    This formulation allows the model to leverage its generative reasoning capabilities to infer sentiment from context, rather than relying solely on pattern-based classification.

    \item \textbf{Low-Rank Adaptation (LoRA):}  
    Fine-tuning a multi-billion-parameter LLM directly is computationally prohibitive. To address this challenge, FinGPT employs Low-Rank Adaptation (LoRA), which freezes the pre-trained model weights and injects trainable low-rank decomposition matrices into selected Transformer layers. This approach drastically reduces the number of trainable parameters (for example, from billions to a few million), enabling efficient adaptation to the linguistic and contextual nuances of the Malaysian financial market without the cost of full-model retraining. Additionally, instruction tuning combined with LoRA improves numerical sensitivity by explicitly training the model to interpret financial quantities within the semantic context of the prompt.
\end{itemize}


\subsection{Prompt Engineering Methodologies}

For the generative component of this study (FinGPT), the methodological focus shifts from traditional parameter optimisation to \emph{prompt engineering}. This approach designs natural language instructions that guide the reasoning process of a Large Language Model (LLM) without necessarily modifying its underlying weights. Unlike FinBERT, which maps textual representations to sentiment classes through a fixed dense classification layer, FinGPT formulates sentiment analysis as a text generation task. As a result, carefully structured prompts are required to align model outputs with the desired sentiment labels (Positive, Negative, Neutral).

\paragraph{Zero-Shot Prompting and Supervised Instruction Tuning}

This research evaluates two complementary generative settings in order to assess baseline generalisation and domain adaptation in the Malaysian market context:

\begin{itemize}
    \item \textbf{Zero-Shot Learning:}  
    In the zero-shot setting, the model is provided only with a task description and a financial news headline, without any labelled examples. The prompt relies entirely on the model's pre-trained knowledge to infer sentiment (e.g.\ \emph{``Return only one label: positive, neutral, or negative''}). This configuration evaluates the model's intrinsic understanding of financial concepts and contextual sentiment.

    \item \textbf{Supervised Instruction Tuning with LoRA:}  
    To capture linguistic, political, and economic nuances specific to Bursa Malaysia, selected generative backbones are further adapted on a supervised headline sentiment corpus using instruction--response pairs and Low-Rank Adaptation (LoRA). This stage evaluates whether gradient-based domain adaptation yields a meaningful improvement over zero-shot prompting alone.
\end{itemize}

\paragraph{Instruction Formatting}

To standardise the generative input, this study adopts the instruction--response format commonly used in instruction tuning. The dataset is restructured into a conversational form:

\begin{quote}
\texttt{Human: [Instruction] + [Input]} \\
\texttt{Assistant: [Output]}
\end{quote}

This structure aligns the sentiment analysis task with the LLM's training objective of next-token prediction, forcing the model to generate a specific sentiment label as a natural language completion. Explicit instructions are particularly important for mitigating numerical sensitivity issues commonly observed in general-purpose LLMs, as they guide the model to interpret quantitative financial statements (e.g.\ \emph{``profit dropped by 5\%''}) within the appropriate semantic and economic context.

\subsection{Evaluation Protocols}

To ensure rigorous and meaningful benchmarking, distinct evaluation metrics are employed for the static classification task (sentiment analysis) and the dynamic regression task (stock price prediction).

\paragraph{Sentiment Benchmarking Metrics (Classification)}

Standard accuracy is often insufficient for financial sentiment analysis due to class imbalance, where neutral news items typically dominate the dataset. To address this limitation, this study employs the following metrics:

\begin{itemize}
    \item \textbf{F1-Score:}  
    The harmonic mean of Precision and Recall, used as the primary evaluation metric for Experiment~1. The F1-score penalises models that achieve artificially high accuracy by predominantly predicting the majority (Neutral) class without detecting market-relevant sentiment shifts:
    \[
        \text{F1} = 2 \times \frac{\text{Precision} \times \text{Recall}}{\text{Precision} + \text{Recall}}.
    \]

    \item \textbf{Precision and Recall:}  
    These metrics are reported separately to diagnose whether a model fails to identify true sentiment signals (low Recall) or produces excessive false positives (low Precision).
\end{itemize}

\paragraph{Stock Prediction Metrics (Regression)}

For the downstream task of predicting KLSE stock price movements in Experiments~2 and~3, model performance is evaluated using error metrics standard in financial econometrics:

\begin{itemize}
    \item \textbf{Root Mean Square Error (RMSE):}  
    RMSE penalises larger prediction errors more heavily than smaller ones, making it particularly suitable for financial forecasting where large deviations carry disproportionate risk:
    \[
        \text{RMSE} =
        \sqrt{
        \frac{1}{n}
        \sum_{i=1}^{n}
        \left( y_i - \hat{y}_i \right)^2
        }.
    \]

    \item \textbf{Mean Absolute Percentage Error (MAPE):}  
    To facilitate comparison across stocks with different price scales (e.g.\ large-cap banking stocks versus low-priced equities), MAPE provides a scale-independent accuracy measure expressed as a percentage:
    \[
        \text{MAPE} =
        \frac{1}{n}
        \sum_{i=1}^{n}
        \left|
        \frac{y_i - \hat{y}_i}{y_i}
        \right|
        \times 100\%.
    \]
\end{itemize}

\paragraph{Validation Strategy: Sliding Window}

To respect the temporal structure of financial time-series data, this study replaces standard $k$-fold cross-validation with a sliding (rolling) window validation strategy. The model is trained on a fixed window of historical data (e.g.\ Days~1--30) and evaluated on the immediately subsequent observation (Day~31). The window then advances forward (Days~2--31) to predict the next step (Day~32). This procedure strictly prevents look-ahead bias, ensuring that future information is never used to predict past events.


\chapter{Sentiment Model Benchmarking}

\chapterabstract{This chapter presents a comparative empirical study of FinBERT and FinGPT-style generative models to evaluate their efficacy in extracting financial sentiment within the Malaysian market context. In addition to the original TinyLlama-based FinGPT baseline, the benchmark is extended to larger open instruct models from the Qwen and Llama families in order to separate the effects of model architecture from backbone capacity. The chapter compares zero-shot prompting, supervised fine-tuning on a Malaysian financial headline corpus, and pseudo-label adaptation on recent unlabeled news. The findings identify FinBERT as the strongest supervised sentiment model for the subsequent forecasting experiments, while also showing that larger generative backbones substantially outperform the original TinyLlama baseline in zero-shot evaluation, whereas public-source FinGPT checkpoints and LoRA fine-tuning of the larger decoders do not surpass either the best zero-shot Llama baseline or FinBERT.}

\section{Introduction}
This chapter evaluates how well domain-specific and open generative language models transfer to Malaysian financial news sentiment classification. The main comparison is between a discriminative encoder-only model, FinBERT, and a set of decoder-only FinGPT-style models that predict sentiment through prompted generation. The motivation is twofold: first, to identify the strongest sentiment model for downstream forecasting; second, to test whether the additional flexibility of generative models yields a practical advantage in an emerging market setting characterised by local political language, corporate reporting conventions, and a strong neutral-class imbalance.

\section{Background Literature}
Prior literature suggests a trade-off between discriminative and generative approaches. FinBERT is computationally efficient and typically strong on conventional benchmark datasets, whereas decoder-only models offer greater contextual flexibility and potentially stronger numerical reasoning. However, the extent to which these strengths transfer to Malaysian news remains an open empirical question. This benchmark therefore treats the generative family not as a single model but as a spectrum of open backbones, ranging from lightweight TinyLlama to larger Qwen and Llama instruct models.

\section{Data Set}
Three data resources were used in the sentiment benchmark.

\begin{itemize}
    \item A supervised headline sentiment corpus of approximately 60,000 labeled financial news headlines, sourced from Kaggle, was used for primary supervised fine-tuning.
    \item A manually annotated Malaysian benchmark set of 100 news items from 2025 was used as the development set for model comparison and selection.
    \item A second manually annotated Malaysian benchmark set of 200 news items from 2026 was reserved as a held-out final test set for the selected downstream model.
\end{itemize}

In addition, a recent Malaysian news corpus covering 2023--2025, with the manually labeled 2025 items excluded, was used for pseudo-label adaptation experiments. This excluded-human corpus contained 92,294 rows after preprocessing.

\section{Aim of Study and Analysis Approach}
The benchmark was structured in three stages.

\begin{enumerate}
    \item \textbf{Zero-shot evaluation.} FinBERT and several open instruct LLM backbones were evaluated directly on the 2025 human-labeled benchmark without task-specific gradient updates.
    \item \textbf{Supervised fine-tuning.} FinBERT and selected generative backbones were adapted on the 60k labeled headline corpus using discriminative fine-tuning for FinBERT and LoRA-based instruction tuning for generative models.
    \item \textbf{Pseudo-label adaptation.} The best 60k-trained FinBERT and TinyLlama-based FinGPT models were used to pseudo-label the 2023--2025 excluded-human corpus, after which each model was adapted on its own pseudo-labeled recent data and reevaluated on the 2025 development benchmark.
\end{enumerate}

Macro F1-score was treated as the primary metric because the Malaysian news benchmark is class-imbalanced and raw accuracy can be inflated by overpredicting the neutral class. Accuracy, macro precision, macro recall, weighted precision, weighted recall, and weighted F1-score were also recorded for diagnostic analysis.

\section{Results}
To keep the chapter aligned with the central thesis question of \emph{FinBERT versus FinGPT}, the main results are organised into two focused comparisons. The first comparison follows the FinBERT family from zero-shot baseline to supervised fine-tuning and pseudo-label adaptation. The second comparison follows the most relevant generative baselines: the official public FinGPT checkpoint, the TinyLlama-based FinGPT run2 model, and the strongest modern open LLM candidate, Meta-Llama-3.1-8B-Instruct. Broader exploratory runs, including Qwen-7B, Qwen-3B, Llama-3.2-1B, Llama-3.2-3B, TinyLlama zero-shot, and the public FinGPT Llama2-13B checkpoint, are treated as supplementary evidence rather than the core narrative.

\begin{table}[htbp]
\centering
\small
\begin{tabular}{|l|l|c|c|}
\hline
\textbf{FinBERT Stage} & \textbf{Model} & \textbf{Accuracy} & \textbf{Macro F1} \\
\hline
Zero-shot baseline & ProsusAI/FinBERT & 0.25 & 0.250 \\
60k supervised & FinBERT run2 & 0.69 & 0.620 \\
Pseudo-label adapted & FinBERT adapted & 0.66 & 0.546 \\
\hline
\end{tabular}
\caption{Main FinBERT-family results on the 2025 human-labeled development set. Macro F1-score is the primary model-selection metric.}
\label{tab:finbert-family-2025}
\end{table}

\begin{table}[htbp]
\centering
\small
\begin{tabular}{|l|l|c|c|}
\hline
\textbf{Generative Stage} & \textbf{Model} & \textbf{Accuracy} & \textbf{Macro F1} \\
\hline
Official public FinGPT & FinGPT MT Llama2-7B (HF) & 0.27 & 0.175 \\
Supervised FinGPT run2 & TinyLlama-based FinGPT run2 & 0.35 & 0.328 \\
Best open zero-shot LLM & Meta-Llama-3.1-8B-Instruct & 0.47 & 0.458 \\
Llama-8B + LoRA (1ep, lr=1e-5) & Meta-Llama-3.1-8B-Instruct + LoRA & 0.44 & 0.431 \\
Llama-8B + LoRA (3ep, lr=2e-6) & Meta-Llama-3.1-8B-Instruct + LoRA & 0.42 & 0.418 \\
\hline
\end{tabular}
\caption{Main FinGPT/LLM comparison on the 2025 human-labeled development set.}
\label{tab:fingpt-family-2025}
\end{table}

Table~\ref{tab:finbert-family-2025} provides the clearest result in the chapter. The supervised FinBERT run2 model dramatically outperformed both the generic zero-shot FinBERT baseline and the adapted FinBERT variant, reaching 0.69 accuracy and macro F1 0.620 on the 2025 benchmark. The pseudo-label adaptation stage remained informative despite being unsuccessful, because it showed that additional recent data did not automatically improve transfer when the labels were noisy. For completeness, the selected 60k-supervised FinBERT model was also evaluated on the held-out 2026 benchmark, where it retained 0.69 accuracy and achieved macro F1 0.598.

Table~\ref{tab:fingpt-family-2025} tells a different story for the generative side. The official public FinGPT checkpoint transferred poorly to the Malaysian benchmark, with macro F1 only 0.175. The TinyLlama-based FinGPT run2 model improved on that baseline and reached macro F1 0.328 after supervised adaptation, but it remained well below FinBERT. The strongest open generative baseline in the focused comparison was Meta-Llama-3.1-8B-Instruct at macro F1 0.458, yet even this stronger decoder model still trailed FinBERT by a wide margin. Moreover, LoRA fine-tuning of Llama-8B consistently degraded performance: the 1-epoch run at lr=1e-5 reduced macro F1 from 0.458 to 0.431, and a further run at 3 epochs with the lower learning rate of 2e-6 reduced it further to 0.418. Reducing the learning rate and increasing the number of epochs did not reverse the decline.

The broader exploratory sweep was still useful, but it is not central to the thesis narrative. Among those supplementary runs, Qwen2.5-7B-Instruct was the second-best zero-shot open LLM at macro F1 0.410, while both Qwen-7B LoRA runs remained below that baseline. The public FinGPT Llama2-13B sentiment adapter also collapsed to near-majority-class behaviour with macro F1 0.133. These additional runs reinforce the main conclusion rather than changing it: larger decoder models improve zero-shot performance, but generative LoRA tuning does not overtake FinBERT in this task.

\section{Summary}
This chapter establishes four focused findings. First, the supervised FinBERT run2 model is the strongest sentiment model in the current Malaysian benchmark and therefore the most appropriate core signal for downstream forecasting. Second, the official public FinGPT baseline is weak in this setting, and the TinyLlama-based FinGPT run2 model improves on it but still remains far below FinBERT. Third, Meta-Llama-3.1-8B is the strongest modern open generative baseline tested in the main comparison, but even it does not match FinBERT, and LoRA fine-tuning consistently degraded performance across all configurations tested (1 epoch at lr=1e-5: macro F1 0.431; 3 epochs at lr=2e-6: macro F1 0.418), confirming that the issue is structural rather than a hyperparameter sensitivity. Fourth, pseudo-label adaptation on recent unlabeled news does not improve performance in this study. The remainder of the thesis therefore treats the 60k-supervised FinBERT model as the primary sentiment extractor, while the generative models serve mainly as comparative evidence on the trade-off between architectural flexibility and classification reliability.




\chapter{Sentiment vs Stock Movement}

\chapterabstract{This chapter empirically investigates the statistical relationship between the aggregated daily sentiment scores derived in Chapter 3 and the market returns of Bursa Malaysia securities under both broad-corpus and company-filtered settings. Given the Chapter~3 benchmark results, the fine-tuned FinBERT signal is treated as the primary sentiment series. The chapter shows that broad market-wide aggregation yields negligible KLSE index signal, whereas company-specific filtering reveals materially stronger same-day associations for selected stocks such as Maybank, CIMB, and Tenaga. These findings reframe sentiment as a relevance-sensitive, primarily contemporaneous signal rather than a universal next-day predictor.}

\section{Introduction}
Chapter~3 identified FinBERT run2 as the strongest sentiment extractor, but classification accuracy alone does not establish whether the resulting signal is economically meaningful. Experiment~2 therefore evaluates whether daily sentiment derived from Malaysian news is associated with observed market movement in Bursa Malaysia. The central question is not only whether sentiment correlates with returns, but also whether the way the news corpus is constructed materially changes the usefulness of the signal.

The experiment was conducted in two stages. First, the full 2023--2025 excluded-human news corpus was treated as a broad market-wide sentiment source and compared with the KLSE index and selected large-cap stocks. Second, the same corpus was filtered into company-specific subsets so that only articles mentioning a given firm contributed to that firm's daily sentiment signal. This design directly tests the hypothesis that weak broad-corpus results may be driven by signal dilution rather than by poor sentiment classification.
\section{Background Literature}
Prior literature on news-driven market prediction suggests that article relevance is a critical determinant of downstream performance. Broad market indices aggregate heterogeneous macroeconomic, political, and sectoral shocks, so a single daily sentiment average may be too diffuse to track any one return series. Event-study and firm-specific sentiment research instead argues that stock-level relationships are strongest when the underlying text is closely tied to the company being evaluated.

This distinction is especially important in the Malaysian setting. Local reporting frequently mixes firm-level announcements with broader political and macroeconomic commentary, and many corporate headlines are neutral or only weakly directional. A model may therefore classify sentiment correctly at the headline level while still generating a weak downstream signal if irrelevant articles are pooled together. Experiment~2 is designed to test exactly this issue by contrasting broad aggregation with company-filtered aggregation.
\section{Data Set}
The primary sentiment source for this chapter was the excluded-human Malaysian news corpus covering 2023--2025, which contained 92,294 preprocessed rows after the manually labeled benchmark items were removed. This corpus was first used directly as a broad market-wide source for daily sentiment aggregation.

To test the effect of article relevance, the same corpus was then filtered into company-specific subsets using ticker and alias matching for leading Bursa Malaysia constituents. The largest filtered corpora used in the detailed analysis were 5,375 Maybank articles (1155.KL), 5,132 CIMB articles (1023.KL), 4,222 Tenaga articles (5347.KL), and 3,246 RHB articles (1066.KL). These sizes were large enough to support daily aggregation across the 2023--2025 period while remaining focused on a single firm.

Market data were obtained from Yahoo Finance for the KLSE index and the selected stocks. Sentiment days were aligned to the same trading day or the next available trading day, and return targets were evaluated at horizons $T$, $T{+}1$, and $T{+}2$.
\section{Analysis Approach}
All sentiment predictions in this chapter were generated with FinBERT run2, the selected model from Chapter~3. For each corpus setting, article-level predictions were aggregated by day and compared with one-day returns at multiple horizons. The analysis reports Pearson correlation, Spearman correlation, directional accuracy, and directional coverage. Macro forecasting models are not yet introduced here; the goal is instead to determine whether the sentiment signal itself has a statistically meaningful relationship with market movement.

The evaluation proceeded in two stages. In the broad-corpus stage, the full market-wide corpus was used to generate a single daily sentiment series and compared against the KLSE index and a wider set of large-cap stocks. In the filtered stage, the corpus was split into company-specific news subsets and each subset was compared only with its corresponding stock. This second stage provides a direct test of whether company relevance strengthens the sentiment--return relationship.
\section{Results}
The broad market-wide analysis produced a largely negative result at the index level. When the full corpus was aggregated into a single daily sentiment series and compared with the KLSE index, correlations were effectively zero at every horizon: Pearson correlation was $-0.008$ at $T$, $-0.001$ at $T{+}1$, and $-0.055$ at $T{+}2$. Directional accuracy also remained close to chance at 0.517, 0.478, and 0.505 respectively. A wider top-10 stock sweep using the same broad corpus produced slightly better directional accuracy for a few individual names, but the correlations remained small; for example, the strongest broad-corpus directional-accuracy values were 0.588 for RHB at $T{+}2$ and 0.568 for Tenaga at $T{+}1$.

It should be noted that the broad KLSE test was conducted over a limited number of matched trading days ($n \approx 31$--49 depending on the lag horizon), which is insufficient for robust statistical inference. The broad-corpus result is therefore treated as indicative rather than conclusive, and the company-filtered analysis, which spans 710--731 matched days per stock, constitutes the primary evidence for this chapter.

These weak broad-corpus results motivated the company-filtered follow-up. Once the news corpus was restricted to company-relevant articles, the same-day relationship became materially stronger for several stocks. Maybank, CIMB, and Tenaga all exhibited moderate positive lag-0 correlation and directional accuracy above 0.64, while RHB improved only modestly. Table~\ref{tab:experiment2-filtered-summary} summarises the contrast between the broad KLSE result and the strongest filtered stock-level results. Table~\ref{tab:experiment2-filtered-lag} presents the full lag profile for each filtered stock.

\begin{table}[htbp]
\centering
\small
\begin{tabular}{|l|l|c|c|c|c|}
\hline
\textbf{Setting} & \textbf{Series} & \textbf{Lag} & \textbf{Pearson} & \textbf{Spearman} & \textbf{Dir. Acc.} \\
\hline
Broad corpus & \texttt{\string^KLSE} & $T$ & -0.008 & 0.039 & 0.517 \\
Broad corpus & \texttt{\string^KLSE} & $T{+}1$ & -0.001 & -0.046 & 0.478 \\
Broad corpus & \texttt{\string^KLSE} & $T{+}2$ & -0.055 & -0.050 & 0.505 \\
Filtered corpus & 1155.KL (Maybank) & $T$ & 0.320 & 0.344 & 0.645 \\
Filtered corpus & 1023.KL (CIMB) & $T$ & 0.342 & 0.365 & 0.645 \\
Filtered corpus & 5347.KL (Tenaga) & $T$ & 0.337 & 0.345 & 0.653 \\
Filtered corpus & 1066.KL (RHB) & $T$ & 0.066 & 0.050 & 0.575 \\
\hline
\end{tabular}
\caption{Comparison between broad-corpus KLSE results and company-filtered stock-level results at lag~0 in Experiment~2. Filtered results are based on 710--731 matched trading days (2023--2025).}
\label{tab:experiment2-filtered-summary}
\end{table}

\begin{table}[htbp]
\centering
\small
\begin{tabular}{|l|c|c|c|c|c|}
\hline
\textbf{Stock} & \textbf{Lag} & \textbf{Pearson} & \textbf{Spearman} & \textbf{Dir. Acc.} & \textbf{Coverage} \\
\hline
1155.KL (Maybank) & $T$     & 0.320  & 0.344  & 0.645 & 0.860 \\
1155.KL (Maybank) & $T{+}1$ & -0.024 & -0.049 & 0.483 & 0.862 \\
1155.KL (Maybank) & $T{+}2$ & 0.002  & 0.011  & 0.502 & 0.863 \\
\hline
1023.KL (CIMB)    & $T$     & 0.342  & 0.365  & 0.645 & 0.885 \\
1023.KL (CIMB)    & $T{+}1$ & 0.052  & 0.031  & 0.518 & 0.885 \\
1023.KL (CIMB)    & $T{+}2$ & 0.044  & 0.020  & 0.521 & 0.886 \\
\hline
5347.KL (Tenaga)  & $T$     & 0.337  & 0.345  & 0.653 & 0.894 \\
5347.KL (Tenaga)  & $T{+}1$ & 0.033  & 0.045  & 0.526 & 0.892 \\
5347.KL (Tenaga)  & $T{+}2$ & 0.087  & 0.088  & 0.531 & 0.894 \\
\hline
1066.KL (RHB)     & $T$     & 0.066  & 0.050  & 0.575 & 0.869 \\
1066.KL (RHB)     & $T{+}1$ & 0.003  & -0.036 & 0.534 & 0.870 \\
1066.KL (RHB)     & $T{+}2$ & 0.008  & -0.009 & 0.531 & 0.866 \\
\hline
\end{tabular}
\caption{Full lag profile for company-filtered stocks in Experiment~2. Coverage denotes the proportion of market days with at least one matched article.}
\label{tab:experiment2-filtered-lag}
\end{table}

Four implications follow from Tables~\ref{tab:experiment2-filtered-summary} and~\ref{tab:experiment2-filtered-lag}. First, the weak KLSE index result is not sufficient evidence that sentiment is useless; rather, it suggests that broad aggregation obscures stock-relevant information. Second, company-specific filtering substantially strengthens the relationship for selected firms, especially Maybank, CIMB, and Tenaga. Third, the effect is sharply contemporaneous: for Maybank, Pearson correlation fell from 0.320 at $T$ to $-0.024$ at $T{+}1$ and near-zero at $T{+}2$, suggesting that any market-relevant information in company-specific news is priced in on the same trading day rather than carried forward. This pattern is consistent with the semi-strong form of the Efficient Market Hypothesis (EMH), which predicts that publicly available news is rapidly incorporated into asset prices, leaving little exploitable lead--lag relationship. Fourth, the weaker result for RHB (Pearson 0.066 at $T$) compared with the other three banks is likely attributable to its smaller filtered corpus (3,246 articles versus 5,132--5,375 for Maybank and CIMB), which reduces the daily sentiment signal-to-noise ratio. The directional coverage column also confirms that approximately 13--14\% of trading days had no matched articles, meaning the sentiment signal is absent for a non-trivial fraction of the evaluation period.

\section{Summary}
This chapter establishes four main findings. First, broad market-wide sentiment derived from the full 2023--2025 corpus has negligible association with KLSE index returns, though this result is based on a limited number of matched trading days and should be interpreted cautiously. Second, company-specific filtering is crucial: once irrelevant articles are removed, FinBERT sentiment exhibits materially stronger same-day relationships for selected firms, most notably Maybank, CIMB, and Tenaga, with Pearson correlations of 0.32--0.34 and directional accuracy above 0.64. Third, the sentiment signal drops sharply beyond lag~0 for most stocks, making it better interpreted as a same-day market-association measure than as a robust next-day forecasting feature; this is consistent with semi-strong market efficiency. Fourth, the signal is not uniform across all firms: RHB's weaker result highlights that corpus size and article relevance quality materially affect downstream performance. These findings motivate the use of relevance-aware, company-filtered sentiment features in the forecasting experiments of Chapter~5, rather than a single broad market-wide aggregate.




\chapter{Stock Prediction Models}

\chapterabstract{This chapter develops and evaluates a hybrid predictive framework by fusing historical market data (OHLCV) with benchmark-selected sentiment features from Chapter~3. Addressing the limitations of traditional technical analysis, the chapter constructs Long Short-Term Memory (LSTM) and Gated Recurrent Unit (GRU) networks to forecast short-term stock prices. The fine-tuned FinBERT model is used as the primary downstream sentiment extractor, while stronger generative baselines remain available as comparative references. The experiment benchmarks the predictive accuracy (RMSE, MAE, MAPE) of sentiment-enhanced models against baseline models relying solely on price history, thereby quantifying the practical forecasting value of localized Malaysian news sentiment.}

\section{Introduction}
\section{Background Literature}
\section{Data Set}
\section{Analysis Approach}
\section{Results}
\section{Summary}




\chapter{Conclusion and Future Work}

\chapterabstract{The final chapter synthesizes the empirical findings regarding the comparative efficacy of encoder-only FinBERT and FinGPT-style generative Large Language Models in Malaysian financial sentiment analysis and forecasting. It summarizes the key contribution of the benchmark study---namely that fine-tuned FinBERT remained the strongest supervised sentiment model, while larger open generative backbones substantially improved the zero-shot baseline relative to TinyLlama. The thesis concludes by identifying limitations in current methodologies and proposing future research directions, including real-time sentiment dashboards, multilingual adaptation, and further evaluation of stronger open generative backbones in downstream market prediction tasks.}

\section{Summary}
\section{Contributions}
\section{Further Work}




% Format a LaTeX bibliography
\makebibliography
\nocite{*}

% Figures and tables, if you decide to leave them to the end
%\input{figure}
%\input{table}

\end{document}

